\documentclass[11pt]{article}

\usepackage[a4paper,margin=29mm]{geometry}
\usepackage[T1]{fontenc}
\usepackage{amsmath,amssymb,amsthm}
\usepackage{booktabs,longtable}
\usepackage[hidelinks]{hyperref}
\hypersetup{
  pdftitle={Marginal-Rank Obstructions to Product-Local Equivalence for Two Pairs of Stabilizer Codes},
  pdfauthor={Xiaojian Zeng}
}

\setlength{\emergencystretch}{2em}
\newtheorem{theorem}{Theorem}
\newtheorem{lemma}{Lemma}
\newtheorem{proposition}{Proposition}
\DeclareMathOperator{\rank}{rank}
\DeclareMathOperator{\Tr}{Tr}
\newcommand{\supp}{\operatorname{supp}}
\newcommand{\cQ}{\mathcal Q}
\newcommand{\cH}{\mathcal H}

\title{Marginal-Rank Obstructions to Product-Local Equivalence\\
for Two Pairs of Stabilizer Codes}
\author{Xiaojian Zeng\\
\small PBC School of Finance, Tsinghua University, Beijing 100190, China\\
\small \texttt{xiaojianzeng@foxmail.com}}
\date{}

\begin{document}
\maketitle

\begin{abstract}
Mahmoud exhibited a five-qubit pair and a six-qubit pair of stabilizer
codes that are related by weight-preserving stabilizer-group isomorphisms
which do not extend to monomial maps, and asked whether the code spaces
might nevertheless be equivalent under arbitrary one-qubit unitaries and a
qubit permutation.  We answer both concrete questions in the negative.
For a normalized stabilizer-code projector $\rho_S$ and a coordinate set
$A$, the marginal rank is
\[
 \rank\!\left(\Tr_{A^c}\rho_S\right)=2^{|A|-d_S(A)},
\]
where $d_S(A)$ is the binary dimension of the stabilizers supported inside
$A$.  This is an explicit rank form of a previously known local-unitary
invariant.  For the five-qubit pair, the two exceptional rank-$4$
three-coordinate sets intersect in one and two coordinates, respectively.
For the six-qubit pair, the rank histograms on four-coordinate sets already
differ.  These permutation obstructions apply to arbitrary product
unitaries, not merely to local Clifford operations.  Complete signed Pauli
and support tables are included, together with a description of exact
reconstruction checks.
\end{abstract}

\section{Context and statement}

The stabilizer formalism describes a quantum code as the common $+1$
eigenspace of an abelian Pauli subgroup not containing $-I$; see, for
example, Gottesman's account~\cite{Gottesman1997}.  Mahmoud's
Theorem~5.3(ii),(iii) gives the following full-support pairs and establishes
their weight-isometry and non-monomial-equivalence properties
\cite{Mahmoud2026}.  Question~5.4 of that paper asks whether either pair is
equivalent at the level of code spaces under a product unitary and a qubit
permutation.

Write $I,X,Y,Z$ for the usual Hermitian one-qubit Pauli matrices.  All
displayed generators below have scalar phase $+1$:
\begin{align*}
 S_A&=\langle ZZZZI,\ XXIIX,\ IIXXX\rangle,&
 S_B&=\langle ZZIZX,\ XIZXI,\ IXXYI\rangle,\\
 S_+&=\langle XXIIXX,\ IIXXXX,\ ZZZZIX\rangle,&
 S_-&=\langle XXZXII,\ ZZXYII,\ IZZZXX\rangle.
\end{align*}
For a stabilizer group $S$, let $\cQ(S)$ denote its common $+1$
eigenspace.  On $n$ qubits, for $\pi\in\operatorname{Sym}(n)$, fix
\[
 P_\pi(v_1\otimes\cdots\otimes v_n)
 =v_{\pi^{-1}(1)}\otimes\cdots\otimes v_{\pi^{-1}(n)}.
\]
Changing this convention to its inverse does not affect the existential
statements.

\begin{theorem}\label{thm:main}
There are no $\pi\in\operatorname{Sym}(5)$ and $U_1,\ldots,U_5\in U(2)$
such that
\[
 (U_1\otimes\cdots\otimes U_5)P_\pi\cQ(S_A)=\cQ(S_B).
\]
Separately, there are no $\pi\in\operatorname{Sym}(6)$ and
$U_1,\ldots,U_6\in U(2)$ such that
\[
 (U_1\otimes\cdots\otimes U_6)P_\pi\cQ(S_+)=\cQ(S_-).
\]
\end{theorem}

The general invariant used below is not introduced here.  Van den Nest,
Dehaene, and De Moor identify the dimensions of the binary stabilizer
subspaces supported in prescribed coordinate sets as degree-two
local-unitary invariants~\cite[Eq.~(14)]{VanDenNestEtAl2004}.  Our
contribution is the explicit application to Mahmoud's two pairs and the
small separating incidence certificates.

\section{Marginal ranks and local-unitary covariance}

Let $[n]=\{1,\ldots,n\}$.  For a Pauli operator $g$, its support is the set
of tensor factors on which it is not the identity.  For $A\subseteq[n]$ put
\[
 S[A]=\{g\in S:\supp(g)\subseteq A\},\qquad
 d_S(A)=\log_2|S[A]|.
\]
For the groups in this note, $d_S(A)$ is equivalently the dimension over
$\mathbb F_2$ of the supported binary stabilizer subspace.

\begin{lemma}[Marginal-rank formula]\label{lem:rank}
Let $S$ be an abelian $n$-qubit Pauli subgroup of order $2^r$ with
$-I\notin S$, and let $\rho_S$ be the trace-one normalization of the
orthogonal projector onto $\cQ(S)$.  Then, for every $A\subseteq[n]$,
\begin{equation}\label{eq:rank}
 \rank\!\left(\Tr_{A^c}\rho_S\right)=2^{|A|-d_S(A)}.
\end{equation}
\end{lemma}

\begin{proof}
The orthogonal code projector and its trace-one normalization are
\[
 \Pi_S=2^{-r}\sum_{g\in S}g,
 \qquad
 \rho_S=\frac{\Pi_S}{2^{n-r}}=2^{-n}\sum_{g\in S}g.
\]
A Pauli term that is nonidentity on a traced-out coordinate has zero partial
trace.  Every term supported in $A$ acquires a factor $2^{n-|A|}$.
Consequently
\begin{equation}\label{eq:partialtrace}
 \Tr_{A^c}\rho_S
 =2^{-|A|}\sum_{g\in S[A]}g|_A.
\end{equation}
The restrictions $g|_A$ form a stabilizer subgroup of order
$2^{d_S(A)}$.  Its normalized group sum is an orthogonal projector of rank
$2^{|A|-d_S(A)}$.  Equation~\eqref{eq:partialtrace} is a nonzero scalar
multiple of that projector, which proves~\eqref{eq:rank}.
\end{proof}

\begin{lemma}[Covariance]\label{lem:covariance}
If two $n$-qubit code spaces are related by a product of one-qubit
unitaries and a coordinate permutation, then their functions
\[
 A\longmapsto\rank\!\left(\Tr_{A^c}\rho\right)
\]
are related by the same permutation of $[n]$.
\end{lemma}

\begin{proof}
Equality of the code spaces gives conjugacy of their uniquely defined
trace-one normalized code projectors.  Under partial trace, unitary factors
on $A^c$ cancel, while those on $A$ conjugate the reduced operator.
Conjugation does not change rank.  A tensor-factor permutation only relabels
the kept set $A$.
\end{proof}

This argument uses no Clifford hypothesis: each $U_i$ may be an arbitrary
member of $U(2)$.

\section{Exact finite certificate}

The three generators in each displayed triple commute and are independent,
so every group below has eight elements.  Table~\ref{tab:paulis} records
all signed products.  If the listed generators are $(g_1,g_2,g_3)$, the
selector $b_3b_2b_1$ denotes $g_1^{b_1}g_2^{b_2}g_3^{b_3}$.  Thus the
table retains the signs forced by Pauli multiplication; they are not
discarded support labels.  A string such as $125$ denotes the subset
$\{1,2,5\}$.

\begin{longtable}{@{}cllc@{}}
\caption{Complete signed Pauli and support tables.}\label{tab:paulis}\\
\toprule
Group & Selector & Signed Pauli word & Support\\
\midrule
\endfirsthead
\toprule
Group & Selector & Signed Pauli word & Support\\
\midrule
\endhead
\midrule
\multicolumn{4}{r}{Continued on next page}\\
\endfoot
\bottomrule
\endlastfoot
$S_A$ & $000$ & $+\mathtt{IIIII}$ & $\varnothing$\\
      & $001$ & $+\mathtt{ZZZZI}$ & $1234$\\
      & $010$ & $+\mathtt{XXIIX}$ & $125$\\
      & $011$ & $-\mathtt{YYZZX}$ & $12345$\\
      & $100$ & $+\mathtt{IIXXX}$ & $345$\\
      & $101$ & $-\mathtt{ZZYYX}$ & $12345$\\
      & $110$ & $+\mathtt{XXXXI}$ & $1234$\\
      & $111$ & $+\mathtt{YYYYI}$ & $1234$\\
$S_B$ & $000$ & $+\mathtt{IIIII}$ & $\varnothing$\\
      & $001$ & $+\mathtt{ZZIZX}$ & $1245$\\
      & $010$ & $+\mathtt{XIZXI}$ & $134$\\
      & $011$ & $-\mathtt{YZZYX}$ & $12345$\\
      & $100$ & $+\mathtt{IXXYI}$ & $234$\\
      & $101$ & $+\mathtt{ZYXXX}$ & $12345$\\
      & $110$ & $-\mathtt{XXYZI}$ & $1234$\\
      & $111$ & $+\mathtt{YYYIX}$ & $1235$\\
$S_+$ & $000$ & $+\mathtt{IIIIII}$ & $\varnothing$\\
      & $001$ & $+\mathtt{XXIIXX}$ & $1256$\\
      & $010$ & $+\mathtt{IIXXXX}$ & $3456$\\
      & $011$ & $+\mathtt{XXXXII}$ & $1234$\\
      & $100$ & $+\mathtt{ZZZZIX}$ & $12346$\\
      & $101$ & $-\mathtt{YYZZXI}$ & $12345$\\
      & $110$ & $-\mathtt{ZZYYXI}$ & $12345$\\
      & $111$ & $+\mathtt{YYYYIX}$ & $12346$\\
$S_-$ & $000$ & $+\mathtt{IIIIII}$ & $\varnothing$\\
      & $001$ & $+\mathtt{XXZXII}$ & $1234$\\
      & $010$ & $+\mathtt{ZZXYII}$ & $1234$\\
      & $011$ & $+\mathtt{YYYZII}$ & $1234$\\
      & $100$ & $+\mathtt{IZZZXX}$ & $23456$\\
      & $101$ & $-\mathtt{XYIYXX}$ & $12456$\\
      & $110$ & $+\mathtt{ZIYXXX}$ & $13456$\\
      & $111$ & $-\mathtt{YXXIXX}$ & $12356$\\
\end{longtable}

\newpage
For completeness, Table~\ref{tab:paulis} and Lemma~\ref{lem:rank} give the
following compact enumeration of every marginal rank, including the empty
and full coordinate sets.

\begin{proposition}[Complete marginal-rank profiles]\label{prop:profiles}
Let $r_S(A)=\rank(\Tr_{A^c}\rho_S)$.
\begin{enumerate}
\item For $S_A$, $r_{S_A}(A)=2^{|A|}$ when $|A|\leq2$.
Among triples the rank is $4$ exactly at $125$ and $345$, and is $8$
otherwise.  Among four-sets the rank is $4$ at $1234$ and $8$ otherwise;
$r_{S_A}([5])=4$.
\item For $S_B$, $r_{S_B}(A)=2^{|A|}$ when $|A|\leq2$.
Among triples the rank is $4$ exactly at $134$ and $234$, and is $8$
otherwise.  Among four-sets the rank is $4$ at $1234$ and $8$ otherwise;
$r_{S_B}([5])=4$.
\item For $S_+$, $r_{S_+}(A)=2^{|A|}$ when $|A|\leq3$.
Among four-sets the rank is $8$ exactly at $1234,1256,3456$, and is $16$
otherwise.  Among five-sets the rank is $8$ at $12345,12346$ and $16$
otherwise; $r_{S_+}([6])=8$.
\item For $S_-$, $r_{S_-}(A)=2^{|A|}$ when $|A|\leq3$.
Among four-sets the rank is $4$ exactly at $1234$, and is $16$ otherwise.
Among five-sets the rank is $8$ at $12345,12346$ and $16$ otherwise;
$r_{S_-}([6])=8$.
\end{enumerate}
\end{proposition}

\begin{proof}
For each $A$, count in Table~\ref{tab:paulis} the group elements whose
support is contained in $A$, take the base-two logarithm, and apply
Lemma~\ref{lem:rank}.  The four stated rules exhaust all $32$ or $64$
subsets.
\end{proof}

\begin{proof}[Proof of Theorem~\ref{thm:main}]
For $S_A$, the three-coordinate sets with marginal rank $4$ are precisely
$125$ and $345$; their intersection is $\{5\}$ and has size one.  For
$S_B$, they are precisely $134$ and $234$; their intersection is
$\{3,4\}$ and has size two.  A coordinate permutation preserves
intersection cardinality.  It therefore cannot identify the two indexed
rank profiles.  Lemma~\ref{lem:covariance} rules out the five-qubit
equivalence.

For $S_+$, the histogram of marginal ranks over the fifteen four-coordinate
sets is
\[
 \{8:3,\ 16:12\}.
\]
For $S_-$, it is
\[
 \{4:1,\ 16:14\}.
\]
A coordinate permutation preserves this unlabelled histogram, so it cannot
identify these two profiles either.  Lemma~\ref{lem:covariance} rules out
the six-qubit equivalence.
\end{proof}

\section{Verification, scope, and disclosure}

The certificate is small enough to check directly: multiply each generator
triple, retain the scalar signs, and read $d_S(A)$ from
Table~\ref{tab:paulis}.  Two accompanying standard-library Python programs
also reconstruct the finite calculation.  The first builds ordinary dense
Pauli matrices over the Gaussian integers, checks commutation, eight
distinct Hermitian products, the projector identity $R_S^2=8R_S$, and the
code rank.  It then forms every partial trace entry by entry and ranks all
$192$ code--subsystem pairs by exact elimination over $\mathbb Q(i)$.  The
second uses binary symplectic labels, independently reconstructs all four
support-dimension profiles, and exhausts all $5!$ and $6!$ coordinate
permutations.  No floating-point tolerance or random choice enters either
check.  From the project root they can be replayed with
\begin{verbatim}
python3 evidence/research_projector_audit.py
python3 evidence/triage_exact_check.py
\end{verbatim}

Separating discovery output from an exact reconstruction program is also
used in finite computational mathematics outside the present setting; one
explicit certificate architecture, based on integer predicates and an
independent implementation, is described in
\cite[Sec.~6]{ZengEtAl2026ExactFinite}.  That paper concerns finite
geometric models and is not an input to the stabilizer argument here.

Theorem~\ref{thm:main} resolves only the two concrete equivalence questions
for the four displayed generator triples, with their specified $+1$
phases.  It does not settle the broader LU--LC-type question in
Mahmoud's Question~5.4.  Agreement of marginal-rank profiles is only a
necessary condition for product-local equivalence; the strict mismatches
above are sufficient for non-equivalence.  Entangling unitaries and
antiunitaries are outside the statement.

OpenAI Codex assisted with primary-source searching, exact-check
organization, and preparation of the manuscript.  Language-model output is
not used as mathematical evidence: the proof rests on the displayed
projector identities and the finite signed table, with the accompanying
exact programs providing reproducible checks.  The author is responsible
for the final claims and presentation.

\bibliographystyle{plain}
\bibliography{references}

\end{document}
